\frontchapter{Executive Summary}
\textbf{Motivation.} Cognitive decline makes everyday recall unreliable:
people misplace objects, forget whether tasks were done, and then rely on
caregivers who often were not present. The loop of manual searching and
caregiver reconstruction (Figure~\ref{fig:recall-loop}) wastes time and effort
for both parties, and the population affected is growing as societies age.

\textbf{Project.} Aegis is an assistive ``external episodic memory''. A
wearable first-person camera observes the user's indoor environment; the system
builds a persistent spatio-temporal memory of objects, layout, and activities,
keeps recent history in fine detail, and condenses older history into
summaries. Natural-language questions (``Where are my keys?'') are answered by
speech or text, with supporting evidence, for the user and, through a separate
caregiver application, for a caregiver (consent-based access is planned). Four
modules, each with a named lead, implement the pipeline
(Section~\ref{sec:modules}).

\textbf{Research findings.} We reviewed three closely related systems --
MemPal, R4DSG, and DAAAM -- plus CrossSense and the wider egocentric-memory
literature (Chapter~\ref{chap:research}). They show that voice-based object
retrieval from a wearable camera is feasible for older adults, that
anchor-relative memory can be built from RGB-only wearable video, and that
hierarchical spatio-temporal memory supports long-horizon questions. Among the
systems reviewed, none combines these with progressive temporal abstraction and
delegated caregiver querying; this is the gap Aegis targets. Since the
mid-evaluation the team ran its own component experiments: a detector
comparison that replaced YOLOE with SAM~3, a memory-compaction audit, and a
seven-day EgoLife study that exposed cross-session room-identity instability
(Section~\ref{sec:experiments}, Chapter~\ref{chap:testing}).

\textbf{Requirements and plan.} Chapter~\ref{chap:requirements} specifies eight
top-level functional requirements with 44 sub-requirements (six added since the
mid report, \tbd{pending supervisor approval}), 16 non-functional requirements,
35 acceptance checks with their FYP-1 results, a traceability matrix and a
requirements change log. Chapter~\ref{chap:plan} records two completed
iterations against plan; Iterations~3 and~4 remain for FYP-2.

\textbf{Implementation and validation status.} Iteration~2 delivered live
capture, a gravity frame, nine spatial relations, state runs, a persistent home
store with cross-session matching, verified compaction, the NLQ engine and
\path{/ask} endpoint on Modal, and patient and caregiver apps
(Chapter~\ref{chap:iter2}). The suite has 451 passing tests at 79\% coverage.
Measured results are detection F1 0.721, compaction 5.7--7.2$\times$ with 100\%
semantic retention, one end-to-end run, and the EgoLife study. Accuracy on a
held-out clip, latency p95, room metrics and usability remain unmeasured.

\textbf{Risks and remaining work.} The three main risks are now: (1)
cross-session room identity, because room linking uses an absolute radius in
scale-free reconstruction units; (2) temporal and activity questions, which the
engine currently declines; and (3) the absence of authentication and
encryption. FYP-2 addresses these in Iterations~3 and~4
(Chapter~\ref{chap:fyp2}).
